\documentclass[11pt,letterpaper]{article}
\usepackage[T1]{fontenc}
\usepackage[utf8]{inputenc}
\usepackage{lmodern,microtype}
\usepackage[margin=0.9in]{geometry}
\usepackage{amsmath,amssymb,amsthm,mathtools}
\usepackage{booktabs,longtable,tabularx,array}
\usepackage{xcolor,enumitem,listings}
\usepackage{hyperref}
\usepackage{fancyhdr}
\definecolor{ink}{HTML}{17324D}
\definecolor{linkblue}{HTML}{245D85}
\hypersetup{colorlinks=true,linkcolor=linkblue,urlcolor=linkblue,citecolor=linkblue,
  pdftitle={Simplifying K215 Proofs: Uniform Decoder Plans and Smaller Proof Dependencies},
  pdfauthor={Research report for Vladimir Reshetnikov}}
\pagestyle{fancy}
\fancyhf{}
\fancyhead[L]{\small\color{ink}K215 proof simplification}
\fancyhead[R]{\small October 21, 2026}
\fancyfoot[C]{\thepage}
\setlength{\headheight}{14pt}
\setlength{\emergencystretch}{2em}
\setlist{nosep,leftmargin=*}
\lstset{basicstyle=\ttfamily\footnotesize,breaklines=true,columns=fullflexible,
  keepspaces=true,frame=single,rulecolor=\color{black!15},backgroundcolor=\color{black!2},
  aboveskip=0.7em,belowskip=0.7em}
\newtheorem{theorem}{Theorem}[section]
\newtheorem{proposition}[theorem]{Proposition}
\newtheorem{lemma}[theorem]{Lemma}
\newcommand{\code}[1]{\texttt{\small #1}}
\newcommand{\Reach}{\mathrel{\longrightarrow^{*}}}
\newcommand{\Reachp}{\mathrel{\longrightarrow^{+}}}
\newcommand{\Con}{\operatorname{Con}}
\newcommand{\Rho}{\rho_8}
\newcommand{\base}{\texttt{2ca3a5c16f0876d2cab842056353c03a1ed2f5f0}}
\newcommand{\repo}{https://github.com/VladimirReshetnikov/ProveIt/blob/2ca3a5c16f0876d2cab842056353c03a1ed2f5f0}
\newcommand{\leanbase}{\repo/Computability/BusyBeaver/Lean/BusyBeaver/ZFC}
\newcommand{\formalbase}{\repo/Computability/BusyBeaver/Research/ZFC/k225-formal}

\title{\color{ink}\textbf{Simplifying the K215 Correctness Proofs}\\[0.4em]
  \large Uniform decoder plans, semantic interfaces, and smaller proof dependencies}
\author{Research report for Vladimir Reshetnikov\\\small ProveIt project}
\date{October 21, 2026}

\begin{document}
\maketitle
\thispagestyle{empty}

\begin{abstract}
The most useful next reductions are structural: expose the finite decoder plan before
quantifying over its tape interpretation; let simulation consume semantic edge contracts;
and separate general set-theoretic lemmas from historical machine endpoints. This packet
supplies a new Lean-checked strengthening of the actual K215 decoder theorem: one fixed
clock and primitive-entry port work for every gap and register list. A checked source-depth
lemma proves that every such prefix makes at least two physical steps. The complete
23-module decoder/runtime dependency closure and both new modules were freshly compiled.
The packet also supplies a focused, independently checked Coq port of Lean's existing
58-step abstract implication proof, two axiom-free Lean enumerator abstractions, and an
uncompiled import-refactoring patch with a reproducible net reduction of 13 local modules
and 233,503 source bytes in the projected K215 dependency closure. These changes preserve
the accepted 215-state machine. Production endpoint integration and the internalization
of its external correctness theorem are explicitly distinguished from the checks completed here.
\end{abstract}

\paragraph{Reviewed baseline.}
All repository-specific findings refer to commit
\begin{center}\small\base\end{center}
The exact dense transition table has
SHA-256
\begin{center}\small\ttfamily
fdbcef0825886035c3091c8537a69535bc96f9e265f0ae45863543d40ae70090
\end{center}
No production repository file was modified during this investigation. The new files and
the proposed patch are delivered for review and incorporation.

\clearpage
\tableofcontents
\clearpage

\section{Recommendation and status}

The final K215 theorem is already short. Its Lean endpoint is a small application of a
shared endpoint theorem; the current enumerator seed theorem is likewise a one-line
instance of a general result. Further gains should come from reducing the proof's
dependencies and eliminating repeated computational work below those interfaces.
Trying merely to compress the final theorem statement would obscure the useful boundaries.
The relevant current sources are \code{KEndpoint215}, \code{SeedPerm}, and the maintained
verification page~\cite{endpoint,verification}.

The recommended order is:
\begin{enumerate}
\item Adopt the stronger uniform decoder plan and its exact zipper adapter. This is a
completed local proof development against the current source and table.
\item Introduce semantic edge coverage, with the current Boolean checker as its first
producer. Then use the shared decoder plus residual tail checks for one existing site batch.
\item Separate general model constructions from the old 254-state endpoint modules.
The provided patch is a concrete source proposal with quantified dependency savings.
\item Give Rocq the same abstract implication interface as Lean, or select the already
available smaller numerical rendering if source size is the sole objective.
\item Share early-halting composition and finite candidate-search lemmas across the
enumerators. The generic lemmas are checked; their concrete instances remain to be connected.
\end{enumerate}

\begin{table}[ht]
\centering\small
\begin{tabularx}{\textwidth}{@{}p{0.25\textwidth}Xp{0.27\textwidth}@{}}
\toprule
Deliverable & Result & Validation in this packet\\
\midrule
Uniform decoder and zipper & One clock and port before all tape/gap/register parameters;
at least two physical prefix steps & Lean 4.32.0; 23 original modules plus 2 new modules\\
Abstract implication core & Same 58-step derivation and interface in Rocq & Coq 8.18.0;
closed reports and independent \code{coqchk}\\
Numerical implication adapter & Preserves \code{DomImp} and \code{rho8\_all} & Focused Coq
8.18 check with documented import-name compatibility change\\
Enumerator proof rules & Composition of partial-result contracts; finite candidate progress &
Two standalone Lean files, ten axiom-free exported results\\
Import refactoring & 180 to 167 BusyBeaver-local modules in projected closure & Source-graph
measurement and patch applicability; integrated build pending\\
Semantic coverage and tails & Reuse decoder proof in production site coverage & Precise
implementation proposal; full site migration pending\\
\bottomrule
\end{tabularx}
\caption{Different validation scopes are deliberately reported separately.}
\end{table}

The report does not claim a new full endpoint build, a smaller transition table, or a
measured production speedup. The new decoder theorem and the standalone proof artifacts
are actual checked results; the integration proposals have their own outstanding obligations.

\section{What the current repository already establishes}

\subsection{The exact theorem and the internalization objective}

The accepted machine has 215 working states, two tape symbols, a two-way infinite tape,
blank input, initial state zero, and a separate halt. Initialization is included in that
state count. The repository records checked external theorems of the form
\[
  \operatorname{Halt}(K215,\mathrm{blank})
  \quad\Longleftrightarrow\quad
  \operatorname{Prov}_{\mathrm{ZFC}}(\bot).
\]
Negating both sides gives nonhalting if and only if syntactic consistency. Consistency
is not assumed in either direction. The repository's verification page explicitly
distinguishes these external proof-assistant results from an internally encoded PA
proof or a derivation in ZFC's own proof syntax~\cite{verification}.

This distinction matters to the stated research objective. An external Lean or Rocq
theorem about a ZFC proof predicate is valuable evidence, but it is not automatically
the separately encoded object
\[
  \mathrm{ZFC}\vdash
  \bigl(\neg\operatorname{Halt}(K215,\mathrm{blank})\leftrightarrow\Con(\mathrm{ZFC})\bigr).
\]
The present work simplifies the available proof architecture. It does not supply that
additional internal derivation. The finite clock construction and explicit derivations
proposed here also make a future arithmetization more transparent, but that is an
architectural benefit, not a completed formalization claim.

\subsection{Avoid repeating completed simplifications}

The reviewed branch already contains three especially relevant improvements.
First, Lean's \code{DominatedImplication.lean} contains a 58-step condensed-detachment
derivation of unrestricted implication row eight, composition, and prefixing. It is
representation independent and its public theorems have no axioms~\cite{rho8}.
Second, \code{KHome} and \code{KHomeReady} already share the scanner's return-to-dispatch
proof. The readiness gate checks fifteen scan rows and an eighteen-step boundary;
the shortcut has positive physical progress~\cite{phase9,home}.
Third, the shared decoder and its exact source and zipper bindings already exist.
They are auxiliary to the endpoint, while the maintained \code{checkEdge} still starts
from the full source-path configuration~\cite{decoder,sim}.

The phase-nine record distinguishes a 192-module endpoint closure from a 197-module
union including five decoder-prefix modules. It records 984.608 seconds for the nine
fresh site batches and 1,649.340 seconds in historical phase-six records. That is a
historical comparison, not a new paired benchmark in this investigation~\cite{phase9}.
The original full build and the later incremental checkpoints also have different
module counts. These numbers must not be mixed with the local import-graph census
in Section~\ref{sec:imports}.

\section{Checked result: choose the decoder plan before interpreting it}
\label{sec:uniform}

\subsection{The unnecessary dependency in the old statement}

The decoder's source word determines a finite route through a role graph. Every edge
on that route names a block with a fixed duration. Its terminal edge determines a
primitive-entry state and the final source bit. These choices depend on the word and
the certified finite graph; they do not depend on tape values outside the word,
the gap length, or the register list.

The existing \code{source\_path\_chain\_routed} nevertheless produces its clock and exit
port inside the universal quantifiers over tape and head position. In simplified form,
its interface is
\[
  \forall a,t,\quad \operatorname{Input}(w,a,t)
       \Longrightarrow \exists N,q,b,\quad \operatorname{Run}(w,N,q,b,a,t).
\]
The later \code{entry\_uniform} selects a port using a particular gap/register instance,
then identifies ports from other instances through route and suffix uniqueness lemmas.
It leaves the exact clock existential inside the parameter quantifiers~\cite{decoder}.

There is no general logical rule that exchanges these quantifiers. The right repair
is a stronger induction that constructs the plan from the source path first.
This is the new theorem in \path{proofs/UniformDecoder.lean}.

\subsection{The strengthened source theorem}

Let \(w\) be a source word, and suppose it describes a valid source-DAG path from the
K215 root to a primitive. Write \(\operatorname{Reads}(t,a,w)\) for the assertion that
the word is present on tape \(t\), starting at integer address \(a\). Write
\(\operatorname{run}_M^N\) for the exact partial execution after \(N\) physical steps.

\begin{theorem}[Uniform K215 decoder plan]
For every such source path, there are natural numbers \(N,q\) and a bit \(b\), chosen
independently of \(a,t\), such that
\[
 |w|-1\le N\le 7(|w|-1),
\]
the corresponding \code{RoutedEntry} certificate holds, and for every head position
\(a\in\mathbb Z\) and full tape \(t\),
\begin{align*}
 t(a-1)=1\ \land\ \operatorname{Reads}(t,a,w)
 \quad\Longrightarrow\quad&
 \operatorname{run}_M^N(84,a,t)
   =\operatorname{some}(q,a+|w|-1,t),\\[-0.3em]
 &t(a+|w|-1)=b.
\end{align*}
The whole tape is preserved. No finite-support hypothesis is imposed on the full tape.
\end{theorem}

\begin{proof}
The proof generalizes over the role identifier and its matching source node, then
inducts on the source path. A source leaf cannot be the current role node because
the role's checked edge specification requires a fork.

For a primitive edge, the path has one remaining bit. Select \(N=0\), the current
role state, and that bit. The route certificate is already finite data. For every
matching tape the chain is reflexive and the last-bit assertion is exactly the
first \code{Reads} hypothesis.

For an internal edge, the checked specification selects a successor role and a block
\(B\). Apply the induction hypothesis to the successor before introducing \(a,t\),
obtaining a fixed tail clock \(N'\), exit state, and final bit. Select
\(N=B.\mathrm{steps}+N'\). Construct the extended route certificate immediately.
Only then introduce the tape. Its \code{Before} and \code{Reads} hypotheses provide
the block window, and the existing framed-block theorem preserves the surrounding
tape. Compose that block with the already selected tail chain.

The bounds follow from \(1\le B.\mathrm{steps}\le7\). The tail word is nonempty,
as witnessed by its terminal-bit split, so natural subtraction in \(|w|-1\) is handled
without an implicit positivity assumption. All choices occur in the finite structural
part of the argument; no choice principle is needed to move the existential witnesses.
\end{proof}

\code{source\_path\_chain\_uniform} is the abstract result.
The theorem \code{reaches\_primitive\_uniform} is the literal K215 instance.
The assumption reports for the generic result
contain \code{propext} and \code{Quot.sound}. The concrete instance also inherits
\code{Classical.choice} from the existing finite-check infrastructure. The claim
that witnesses are selected structurally is therefore not a claim that every imported
theorem in the concrete dependency closure is axiom-free.

\subsection{Checked positivity from three levels of the source}

The literal source root is node 201, with children 16 and 200. Those have children
\((2,15)\) and \((32,199)\), respectively; all four of those nodes are forks.
Consequently every complete source word has at least three bits.

The new file defines a small Boolean \code{atLeastDepth} inspection and proves its
generic soundness by induction: depth zero is vacuous; positive depth excludes leaves
and decreases through the selected fork child. The exact depth-three fact is checked
by \code{decide +kernel}. Therefore
\[
 \operatorname{SourcePath}(w)\Longrightarrow 3\le |w|
 \Longrightarrow 2\le N.
\]
The positivity theorem requires no extra caller-supplied depth premise. This is a
kernel-checked statement about every source path, not an inference from a finite run
sample or a trusted assertion supplied by Python.

\subsection{The exact zipper theorem}

The second new file, \path{proofs/UniformZipper.lean}, connects the uniform plan to
the repository's existing zipper semantics and literal expanded source. Its result is
\begin{align*}
 \exists\,&\mathrm{pre},b,q,\mathrm{op},\mathrm{arg},N:\quad
 w=\mathrm{pre}\mathbin{+\!+}[b],\quad
 e=\mathrm{primitive}(\mathrm{op},\mathrm{arg}),\\
 &\operatorname{RoutedEntry}(w,q,b,\mathrm{op},\mathrm{arg}),\quad
 2\le N,\quad |\mathrm{pre}|\le N\le7|\mathrm{pre}|,\\
 &\forall n,\mathrm{vs},\quad
 \operatorname{run}_M^N(\operatorname{conc}(n,\mathrm{vs},\operatorname{start}(w)))
 =\operatorname{some}\bigl(\operatorname{conc}(n,\mathrm{vs},
              \operatorname{leafPort}(\mathrm{pre},b,q))\bigr).
\end{align*}
The clock is outside \(\forall n,\mathrm{vs}\). The theorem has no register-validity
or nonempty-register-list premise because it stops before the primitive action.
Its proof uses the existing whole-tape-to-zipper reflection and preserves the complete
reversed prefix, marker, gap, and register frame.

\code{walk\_prefix\_uniform} supplies the same result for the actual interpreter walk,
using \code{consume P d}. This retains implicit zero padding; substituting
\code{P.take d} would change the contract. \code{prefix\_tail\_sound} then composes
the positive prefix with any valid ordinary-reachability tail. The old uniqueness
lemmas are not needed by this new construction. Removing them from production is a
later API decision: existing external callers must be checked, and compatibility
wrappers may be preferable.

\subsection{What was compiled}

A fresh serial Lean 4.32.0 build compiled the 23 original modules reachable from
\code{KDecoderPrefix215}, followed by the two new modules. All 25 succeeded. The
actual literal decoder, table checks, source binding, symbolic runtime, and old
simulation theorem were compiled in that closure; no placeholder model was substituted.
The final logs and per-source SHA-256 values are in \path{evidence/lean-uniform/}.

The new leaf and walk theorems report only the same three standard Lean axioms:
\code{propext}, \code{Classical.choice}, and \code{Quot.sound}. There is no
\code{sorryAx}. This is a local checked extension of the production semantics, not
a new full build of the ZFC endpoint or a completed adoption of the new site producer.

\section{Use the shared prefix through semantic coverage}
\label{sec:semantic}

\subsection{The interface that currently prevents reuse}

The current \code{KSim.Covered M p} says that every walk reaches a leaf whose
\code{checkEdge M Q e} is true. That Boolean checker executes from
\code{start Q} through the whole path and primitive. The simulation proof repeatedly
unpacks the Boolean expression to recover the semantic facts it actually needs:
path length, register bounds, branch preconditions, exact register effects, and
positive reachability~\cite{sim}.

Define a semantic proposition \code{EdgeRealizes M Q e} containing those facts
directly. Its length field is the current \(|Q|\le48\) guard. Increment and decrement
also retain \(r<9\). Each live branch universally quantifies the natural gap and
register list, retains its \code{Valid} premise, and reaches the identical
\code{tgtConf} expression currently used by \code{checkEdge}. The halt branch asserts
actual halting, and a \code{node} edge is impossible.

For example, the increment case requires, for every valid register frame,
\[
 \operatorname{conc}(n,\mathrm{vs},\operatorname{start}(Q))
 \Reachp
 \operatorname{conc}\bigl(n,\mathrm{vs},
  \operatorname{tgtConf}(\operatorname{incPath}(Q),\operatorname{allOnes}(Q),
       \operatorname{inc}(\operatorname{physOf}(r)))\bigr).
\]
The decrement case includes both the zero and nonzero branch contracts. Break and
continue retain their existing ancestor-path transformations, overflow conditions,
and targets. A semantic interface must not silently discard any of those distinctions.

Then set
\[
\operatorname{SemanticCovered}(M,p):=
 \forall P,\ \exists e,d,\quad
 \operatorname{walk}(p,P)=\operatorname{some}(e,d)
 \ \land\ \operatorname{EdgeRealizes}(M,\operatorname{consume}(P,d),e).
\]
The omitted walk arguments here are exactly the current node bound, root, and initial
depth; the implementation should preserve them literally.

\subsection{A conservative migration}

First prove \code{checkEdge\_sound} using the existing \code{checkEdge\_length},
\code{runTo\_sound}, and \code{runHalt\_sound}. The existing Boolean \code{Covered}
then implies \code{SemanticCovered}. Refactor \code{sim\_step} and \code{sim\_halt}
to use that interface. The register arithmetic, spare-list extension, path cleanliness,
and overflow proofs are retained. \code{sim\_stop} still requires total walk coverage
and the exclusion of node leaves. Preserve the public \code{halts\_iff} as a wrapper.
At this stage all existing clients can continue using exactly the old certificates.

Next add a second producer: use the shared prefix theorem to reach \code{leafPort},
then run a residual checker to the unchanged full target. Initially retain every
source occurrence and both decrement outcomes. This is a smaller step than proving
all parameterized primitive families at once. A mixed rollout can choose the old
full check or the new prefix-plus-tail proof for each occurrence.

For the residual checker, initialize the executor's \code{moved} flag to true.
Its soundness provides ordinary reachability from the primitive-entry port; the
prefix supplies the needed physical progress:
\[
  s\Reachp p\quad\land\quad p\Reach t
       \quad\Longrightarrow\quad s\Reachp t.
\]
The existing \code{Reach1.trans\_reach} proves the composition, and the new
\code{prefix\_tail\_sound} specializes it to the uniform decoder contract. A zero-length
tail is harmless here. This argument depends on positive prefix progress; it must not
be exported to a source with a one-bit leaf without an explicit progress premise.

\subsection{Why the converse halting implication is preserved}

One must preserve positive \emph{physical} progress for every live source step. In the
existing deterministic simulation, if the source takes a step, the target makes at
least one step and reaches a related configuration. A finite halting target run
therefore has a strictly smaller remaining time after a simulated source step.
Strong induction on that time yields the reverse halting implication. A proposed
macrostep that could silently loop through zero-duration administrative work would
not support this induction.

The new interface does not alter that proof. It changes the producer of its premises.
The full prefix retained by \code{leafPort} also matters: two predecessor bits suffice
for the checked internal decoder windows, but they do not summarize every primitive
return context. Break and continue depend on ancestor boundaries; overflow can move
the marker. Generalizing primitive families is appropriate only after their complete
frame contracts have been proved.

\subsection{Reproducible decoder work census}

\path{code/decoder_prefix_census.py} independently parses the pinned decoder and
transition table, checks all 281 local block traces, checks primitive transition
agreement, traverses every source/role occurrence, and literally simulates each full
decoder prefix. It rejects reads outside the specified finite window and requires
restoration of every cell. It imports no repository generator. The result and source
hashes are in \path{evidence/decoder_prefix_metrics.json}.

\begin{table}[ht]
\centering\small
\begin{tabular}{@{}lr@{}}
\toprule
Quantity & Measured value\\
\midrule
Expanded source leaves & 4,751\\
Branch outcomes, counting both decrement cases & 6,176\\
Source DAG nodes / decoder roles / distinct blocks & 202 / 234 / 281\\
Exposed primitive-role interfaces & 97\\
Source-word lengths & 3--47\\
Exact decoder-prefix clocks & 2--55\\
Sum of prefix clocks, once per leaf & 165,876\\
Sum of prefix clocks, once per branch outcome & 213,945\\
Distinct local block traces, checked once each & 353 physical steps\\
Internal edge occurrences / their trace-step sum & 4,749 / 5,721\\
\bottomrule
\end{tabular}
\caption{Finite decoder-only evidence. The clocks and counts do not predict kernel time.}
\end{table}

The 213,945 figure quantifies physical prefix work that is repeated in a full
branch-by-branch interpretation. It is not a symbolic-executor instruction count,
kernel reduction count, or estimated speedup percentage. The minimum depth and
positive progress also have the separate kernel proof in Section~\ref{sec:uniform}.
The exact maximum of 55 in this table remains finite census evidence; the checked
uniform theorem states the structural upper bound \(7(|w|-1)\).

The largest repeated primitive families are increment 1 at state 57/read 0 and
continue 2 at state 57/read 1, each with 267 occurrences; the analogous increment
6 and continue 3 families at state 62 each occur 257 times. These are useful later
targets. The already designed 86-occurrence CB1 case may still be the easiest first
parametric pilot. Occurrence count alone does not establish a correct generic
primitive theorem or determine its proof difficulty.

\section{Separate model theory from historical program endpoints}
\label{sec:imports}

\subsection{Three overly broad dependency routes}

The current K215 import closure contains the historical 254-state enumerator and
register-program proof even though the K215 endpoint consumes more general semantic
results. The key paths include
\begin{align*}
\code{ConstructibleChoice}&\to\code{ZFNative}\to\code{NativeTheory254}\to\code{L2Root254},\\
\code{KQ1ZF}&\to\code{WFTheory}\to\code{WFModel}\to\code{QuotientZF0}\to\code{Enum254},\\
\code{L2Lib222}&\to\code{Enum254}.
\end{align*}
The first consumer needs \code{ZFNative.models\_ZF}. The second needs the generic
well-founded-part model construction, especially \code{WFModel.domain\_models\_ZF}.
The third uses the underlying numerical code layer. These module dependencies are
wider than those uses~\cite{modelimports}.

The supplied patch moves \code{models\_ZF} verbatim to \code{ZFModels.lean} and moves
the generic prefix of \code{WFModel.lean} verbatim to \code{WFModelCore.lean}. The
original files import the new cores and retain their historical specializations,
with declaration names preserved. Consumers import the narrower modules.
\code{KQ1ZF} explicitly imports \code{InductiveAttempts}, retaining the infinity
lemma that previously arrived transitively. \code{L2Lib222} imports \code{Code254}
directly; its maintained generator receives the matching substitution.

This is not deletion of the old mathematics. It keeps those theorems available for
their original callers while removing them from an unrelated endpoint's build chain.
The proposed files, patch, and generator are provided under \path{patches/} and
\path{code/import_boundary_proposal.py}.

\subsection{Exact scope of the projected reduction}

The census counts modules physically under \path{Computability/BusyBeaver/Lean} and
treats seven external library names as boundary nodes. The boundary set is unchanged.
It does not claim to reproduce the verification page's 192-module project-wide
closure count.

\begin{table}[ht]
\centering\small
\begin{tabular}{@{}lrrr@{}}
\toprule
Scoped endpoint closure & Baseline & Proposed graph & Net reduction\\
\midrule
BusyBeaver-local modules & 180 & 167 & 13\\
Source lines & 36,261 & 32,296 & 3,965\\
UTF-8 source bytes & 1,989,076 & 1,755,573 & 233,503\\
\bottomrule
\end{tabular}
\caption{Static source-graph projection, including the new core modules.}
\end{table}

Fifteen historical modules leave the closure and two core modules enter it. The net
byte reduction is approximately 11.74\% in this measured scope. These are source-size
and dependency measurements, not proof-term or elapsed-time measurements.

The patch passes \code{git apply --check}. The generator checks byte-identical
extraction of the moved declaration regions and records source hashes. The integrated
patch has not been compiled. In Lean, narrowing imports can remove unqualified names,
instances, attributes, or simplification lemmas; a textual use analysis cannot settle
every elaboration dependency. Acceptance therefore requires an affected-module build,
the two unchanged K215 endpoint clients, and the legacy clients of the compatibility
wrappers. This is a bounded, concrete proposal with a substantial potential benefit,
not a replacement for that kernel check.

\section{One abstract implication proof in both assistants}

\subsection{What to share}

The numerical row is
\[
 \Rho(a,b,c,u)=((c\Rightarrow u)\Rightarrow a)
       \Rightarrow((a\Rightarrow c)\Rightarrow(b\Rightarrow c)).
\]
The dominated calculus initially grants it only under \(u\preccurlyeq a\).
Lean already derives the unrestricted row using a representation-independent
58-step proof. The production Rocq \code{DomRho8.v} at this baseline still contains
the 22,673-byte compact certificate replay, while \code{rho8cd} also contains an
11,190-byte, 52-step numerical rendering proposal~\cite{rho8,rocqproduction}.

The new \path{rocq/DominatedImplication.v} ports the current abstract Lean derivation
to Rocq. Its interface needs only an arbitrary carrier \(A\), a binary operation
\(\Rightarrow\), a relation \(\preccurlyeq\), and a predicate \(T\), with:
\begin{align*}
 &x\preccurlyeq x,\\
 &u\preccurlyeq x\Longrightarrow u\preccurlyeq(x\Rightarrow y),\qquad
 u\preccurlyeq y\Longrightarrow u\preccurlyeq(x\Rightarrow y),\\
 &T(a\Rightarrow b)\land T(a)\Longrightarrow T(b),\\
 &u\preccurlyeq a\Longrightarrow T(\Rho(a,b,c,u)).
\end{align*}
It proves unrestricted \code{rho8\_all}, composition, and right prefixing. It does
not assume transitivity or antisymmetry of the relation, injectivity of the operation,
a formula decoder, arithmetic, a truth constant, or a logical interpretation of codes.
In particular it never needs the numerical equation \code{wim 0 0 = 0}.

The numerical adapter preserves the existing \code{DomImp} record, including its
truth field, and the exact public \code{rho8\_all} statement. It instantiates the
growth laws using the established pairing bounds. Preserving the unused record field
keeps the interface compatible while allowing the mathematical core to remain smaller.

\subsection{Why propositional automation is not a substitute}

Here \(\Rightarrow\) is an operation on codes and \(T\) is an arbitrary derivability
predicate. A propositional tautology prover would prove something about logical
implication, not about membership in \(T\). Identifying those two levels would require
another soundness or interpretation theorem. Explicit applications of modus ponens
avoid this obligation and remain valid for noninjective encodings and junk codes.

The 58 ordinary proof lemmas apply earlier results at explicit substitution instances.
There are 33 discharged guard occurrences: 15 reflexive, 5 supplied by local hypotheses,
and 13 established by short left/right operand paths. The latter paths are
\(R\) nine times, \(RR\) twice, \(LL\) once, and \(LLR\) once. Both growth laws are
used; these counts do not justify deleting either law.

\subsection{Validation and the real tradeoff}

The port and numerical adapter were compiled with Coq 8.18.0. The generic results
and numerical \code{rho8\_all} print closed assumption reports, and independent
\code{coqchk} checks returned exit zero. For the numerical adapter, the sole
version-compatibility source change was \code{From Stdlib} to \code{From Coq} in
the arithmetic import. The two archived numerical dependencies are hash pinned.
The modern source, compatibility patch, and receipts are supplied separately.

This is not a production Rocq 9.1.0 endpoint build. Replaying the current exact Lean
abstract core under Lean 4.32.0 also succeeded, with all three public results axiom-free.
The mathematical derivation is inherited from the repository; the contribution here
is the common Rocq abstraction and an independent explicit-instance audit.

\begin{table}[ht]
\centering\small
\begin{tabular}{@{}lr@{}}
\toprule
Rocq proof source & Bytes\\
\midrule
Current production replay & 22,673\\
New generic core & 14,655\\
New numerical adapter & 1,389\\
Generic core plus adapter & 16,044\\
Existing numerical 52-step proposal & 11,190\\
\bottomrule
\end{tabular}
\caption{Two different choices: the common abstract interface or the smallest numerical file.}
\end{table}

The proposed shared version saves 6,629 bytes, or approximately 29.24\%, against the
current production replay. It is larger than the existing numerical 52-step proposal.
Its advantage is sharing the same mathematical interface and derivation with Lean.
The row-only dependency slice of the abstract derivation contains 52 lemmas; the six
additional helpers serve composition and prefixing. Whether to retain the full common
interface or specialize the Rocq consumer should be an explicit maintenance choice.

The included \path{rocq/check_and_render.py} checks an expanded certificate containing
explicit substitution arguments. It performs no unification or proof search: it checks
backward references, substitution scopes, exact MP equations, guard witnesses, and
the three intended public goals. Five malformed controls are rejected. The generated
ordinary proof is subsequently checked by the kernel, so the Python checker is not
part of the logical trusted base. Its JSON and code are review artifacts and are not
counted as proof-source savings.

\section{Reusable enumerator contracts and finite progress}

\subsection{Early halt is an allowed result, with a separate soundness invariant}

The current command-block contract says that a finite block either halts early or
returns a specified state. Abstract its repeated case split as
\[
 \operatorname{Returns}(r,x):\Longleftrightarrow
       r=\operatorname{none}\ \lor\ r=\operatorname{some}(x).
\]
For option composition, the two basic rules are
\begin{align*}
 \operatorname{Returns}(r,x)\land\operatorname{Returns}(f(x),y)
   &\Longrightarrow\operatorname{Returns}(r\mathbin{\mathrm{bind}}f,y),\\
 \operatorname{Returns}(r,x)\land f(x)=\operatorname{none}
   &\Longrightarrow r\mathbin{\mathrm{bind}}f=\operatorname{none}.
\end{align*}
Both follow by the two possible results of \(r\). List concatenation versions follow
from \code{exec\_append}. The included \path{proofs/OptionContracts.lean} proves these
rules and prefix variants using only \code{Init}. All eight printed exported
dependency reports are axiom-free under Lean 4.32.0.

The intended application is to \code{gen\_contract}, \code{mp\_contract}, and
\code{program\_none}: rewrite the existing execution through \code{cexec\_append},
apply the generic composition rule, and prove only the concrete suffix behavior.
For modus ponens, apply the two-prefix rule to its two compiled subproofs. The current
\code{EnumPerm} proofs are already short; the principal benefit is sharing the pattern
with the older and future command encodings~\cite{enumerator}.

In this application \code{none} represents an actual halting event reached during
finite execution. It does not mean divergence. The permissive contract alone is not
a soundness theorem: an erroneous block that always halts satisfies it. The global
invariant proving that every halt yields a refutation must remain in the proof.

\subsection{A scheduler theorem without an unexplained fairness assumption}

Suppose \(\operatorname{Ready}(s)\) means that execution is between candidates, and
\(\operatorname{next}(s)\) is the next candidate counter. The concrete fetch/drain
proof should establish
\[
 \operatorname{Ready}(s)\Longrightarrow
 \left(\exists t,\ s\Reach t\land\operatorname{Ready}(t)
          \land\operatorname{next}(t)=\operatorname{next}(s)+1\right)
 \ \lor\ \operatorname{Halts}(s).
\]
Finite reachability must be reflexive and transitive, and halting must propagate
backward along it. For an initial state \(s\) satisfying \(\operatorname{Ready}(s)\),
induction on \(k\) then proves that counter
\(\operatorname{next}(s)+k\) is finitely reached or the original state already halts.

If a compiled reset-and-refute block has a sufficiently late candidate code, its
contract provides a \(k\) for which every ready state at that counter halts. Compose
that fact with the finite-progress theorem to obtain halting from the original state.
\path{proofs/FairSearch.lean} proves both statements with no axioms and only
\code{Init}. Its name does not hide an infinite fairness assumption: all progress
premises are finite and must be discharged by the concrete enumerator.

The fetch, drain, command encoding, stack invariants, and killer-candidate construction
remain substantive concrete obligations. The generic theorem shares their assembly;
it does not make those obligations assumptions of the final machine endpoint.

\section{Integration plan and acceptance boundaries}

\paragraph{First change: uniform decoder information.}
Incorporate the stronger source-path theorem and literal depth-three fact. Replace
the later uniformity reconstruction by the checked zipper adapter. Keep older public
statements as consequences if they have callers. This makes progress and clock
uniformity explicit before changing the production site-checking strategy.

\paragraph{Second change: semantic coverage.}
Introduce the semantic edge record and an adapter from every existing Boolean site
check. Rebuild the unchanged simulation consumers before introducing a new producer.
Then route one current site batch through exact prefix plus residual tail, retaining
the same table, targets, branch cases, and register preconditions. Compare this batch
against the old producer on the same toolchain and machine. Only after that comparison
should a production speedup be stated or the remaining batches migrated.

\paragraph{Third change: narrower module imports.}
Apply the supplied patch in an isolated checkout. Rebuild the changed model cores and
their descendants, then check both endpoint clients and the historical wrappers.
Regenerate the L2 source to verify that the import improvement survives its producer.
Refresh the source-bound verification record after these builds succeed.

\paragraph{Fourth change: common Rocq implication core.}
Compile the generic file and the numerical adapter with the repository's production
Rocq 9.1.0 and stdlib configuration. Update the generator, module order, and manifest
to reproduce the new files. Compile \code{ThmDom} and its actual dependents, then run
the recursive endpoint check before declaring the new lane adopted. The current
checked replay remains a reproducible baseline during this process.

\paragraph{Fifth change: enumerator instances.}
Connect the option-contract rules to the existing finite execution semantics and
instantiate the scheduler from the current fetch/drain lemmas. The existing seed
theorem already has a small public wrapper; aim to reduce duplicated compiler and
scheduler reasoning rather than changing that wrapper's statement.

Each change has a separate acceptance point. This is particularly useful here because
the repository already records merges between independently developed proof lanes.
A successful historical check belongs to its exact source bytes and imported
definitions. A new producer, relocation, or executable-definition change should not
inherit an old check merely because a final theorem statement is unchanged.

\section{Archive contents and reproduction}

\begin{longtable}{@{}p{0.34\textwidth}p{0.61\textwidth}@{}}
\toprule
Path & Purpose\\
\midrule
\endhead
\path{report.tex}, \path{report.pdf} & Complete report and compiled document.\\
\path{proofs/UniformDecoder.lean} & Strong source-plan theorem; depth check and positivity.\\
\path{proofs/UniformZipper.lean} & Literal source leaf/walk adapters and prefix-tail composition.\\
\path{proofs/OptionContracts.lean}, \path{proofs/FairSearch.lean} & Standalone checked enumerator proof abstractions.\\
\path{rocq/} & Generic implication core, numerical adapter, explicit certificate,
checker, compatibility patch, and focused receipts.\\
\path{patches/import_boundaries.patch} & Reviewable import refactoring; integrated compilation pending.\\
\path{patches/proposed/} & Proposed new/edited files for inspection.\\
\path{code/} & Reproduction of the decoder checks, proof build, and import census.\\
\path{evidence/lean-uniform/} & Final 25-module source-bound build record and raw outputs.\\
\path{evidence/} & Decoder census, import projection, standalone proof logs, and source review.\\
\path{README.md}, \path{REPRODUCE.md} & Status map, build commands, and remaining integration obligations.\\
\path{NOTICE.md}, \path{LICENSE} & Attribution and inherited license boundary.\\
\path{MANIFEST.sha256} & Artifact integrity inventory, excluding itself.\\
\bottomrule
\end{longtable}

The Lean decoder build takes a checkout at the pinned commit and a Lean 4.32.0
executable. It discovers and freshly compiles the 23 original dependencies, then the
two new files, recording each source hash and compiler exit status. It needs only the
bundled Lean library for this local closure, not a full mathlib or ZFC-endpoint build.
The standalone enumerator lemmas likewise require only \code{Init}.

The Coq core uses the default prelude. Its numerical adapter pilot uses two pinned
arithmetic source modules archived in the repository; the reproduction script applies
the explicit import-name compatibility change for Coq 8.18. The production Rocq 9.1
adoption should use the modern adapter source and its actual maintained dependencies.

All source/proof artifacts are text. Compiled Lean and Coq objects and downloaded
toolchains are omitted from the archive. The final build records provide their hashes,
commands, logs, and validation scope. The report's recommendation is therefore concrete:
adopt the newly checked uniform decoder information first, expose semantic coverage
so it can be used, and reduce unrelated dependencies without changing the accepted
machine or weakening its correctness statement.

\begin{thebibliography}{12}
\small
\bibitem{verification}
ProveIt, \emph{Formal correctness of the exact 215-state machine}, pinned verification page.
\href{\formalbase/verification/k215/README.md}{Repository source and recorded assurance boundary}.

\bibitem{endpoint}
ProveIt, \href{\leanbase/KEndpoint215.lean}{\code{KEndpoint215.lean}} and
\href{\leanbase/SeedPerm.lean}{\code{SeedPerm.lean}}, pinned sources.

\bibitem{phase9}
ProveIt, \emph{Ninth K215 proof-simplification checkpoint}.
\href{\formalbase/verification/k215/simplification/phase9/README.md}{Phase-nine status and timing evidence}.

\bibitem{home}
ProveIt, \href{\leanbase/KHome.lean}{\code{KHome.lean}} and
\href{\leanbase/KHomeReady.lean}{\code{KHomeReady.lean}}, shared return proof and gate.

\bibitem{decoder}
ProveIt, \href{\leanbase/DecoderPaths.lean}{\code{DecoderPaths.lean}},
\href{\leanbase/DecoderSourceBinding.lean}{\code{DecoderSourceBinding.lean}},
\href{\leanbase/KDecoderPrefix215.lean}{\code{KDecoderPrefix215.lean}}, and
\href{\leanbase/K215Decoder.lean}{\code{K215Decoder.lean}}, pinned decoder sources.

\bibitem{sim}
ProveIt, \href{\leanbase/KSim.lean}{\code{KSim.lean}} and
\href{\leanbase/KRuntimeCore.lean}{\code{KRuntimeCore.lean}}, coverage and simulation.

\bibitem{rho8}
ProveIt, \href{\leanbase/DominatedImplication.lean}{\code{DominatedImplication.lean}},
\href{\leanbase/DomRho8.lean}{\code{DomRho8.lean}}, and
\href{\formalbase/rho8cd/README.md}{\code{rho8cd} derivation and producer notes}.
The implication derivation is inherited mathematical work.

\bibitem{rocqproduction}
ProveIt, \href{\formalbase/rocq/rocq-extension/coq/215/DomRho8.v}{production \code{DomRho8.v}},
\href{\formalbase/rocq/rocq-extension/coq/215/Source215.v}{\code{Source215.v}}, and
\href{\formalbase/rho8cd/PROVEIT_STATUS.md}{lane-specific adoption status}.

\bibitem{modelimports}
ProveIt, \href{\leanbase/ZFNative.lean}{\code{ZFNative.lean}},
\href{\leanbase/WFModel.lean}{\code{WFModel.lean}},
\href{\leanbase/KQ1ZF.lean}{\code{KQ1ZF.lean}},
\href{\leanbase/ConstructibleChoice.lean}{\code{ConstructibleChoice.lean}}, and
\href{\leanbase/L2Lib222.lean}{\code{L2Lib222.lean}}, source of the import projection.

\bibitem{enumerator}
ProveIt, \href{\leanbase/EnumPerm.lean}{\code{EnumPerm.lean}} and
\href{\leanbase/EnumDom.lean}{\code{EnumDom.lean}}, compiler contracts and candidate enumeration.

\bibitem{attribution}
ProveIt, \href{\repo/Computability/BusyBeaver/Research/ZFC/candidate-k215/NOTICE.md}{K215 attribution notice}.
Inherited mathematics and machine ancestry remain credited to their original contributors,
including the ProveIt contributors, Elliot Glazer, Jeff Ketchersid, and the other authors
identified in those notices.
\end{thebibliography}

\end{document}
